As demonstrated by the tests performed and discussed in the previous chapter, the current implementation 
has some limitations.
In this chapter, I will discuss the main limitations of the current implementation and try to outline 
a direction for future work in this field.

\section{Limitations}
    The main limitation is that the developed tool fails to correctly modify some binaries, as 
    shown in \cref{tab:experiment1}.
    I have identified some possible causes of instrumentation failure.
    One of these occurs in some binaries, where the addition of segments (needed to hold the new sections) 
    causes other sections in other segments to shift.
    In fact, according to the \textbf{LIEF} documentation, LIEF does not guarantee preservation of the \gls{ELF} layout.
    For instance, when the dynamic dependency is added to the binary, some sections are modified. Specifically:
\begin{itemize}
\item \textit{.dynstr}: the string 'libforkserver.so' must be added to the dynamic string section.
\item \textit{.dynamic}: LIEF adds a DT\_NEEDED entry that points to the newly added string in \textit{.dynstr}.
\end{itemize}
    If, during either of these two additions, LIEF exceeds the available space in the segment containing the 
    grown section, it reorganizes the segments.
    This causes a complete displacement of the data sections (such as \textit{.bss, .data, .rodata}, \dots) 
    without a corresponding fix to the instructions that reference these data sections.
    This can cause a segmentation fault when, for example, an instruction tries to access a value stored 
    in \textit{.bss}.
    In both PIE and non-PIE binaries, the non-linear displacement between the \textit{.bss} section's address 
    before and after the binary modification causes instructions to no longer point to the correct address.
    This fact can be seen as the underlying problem both for \textit{readelf} and \textit{sfconvert} of \cref{tab:experiment1}. 
    In \textit{who} of \cref{table:fuzzingStats} instead, the problem needs to be traced in some error inside the instrumentation 
    that cause the disalignment of the instructions of a \gls{BB}, causing some of these instruction to be illegal instructions.\\ 
    Moreover, for other binaries, the problem arise in the actual instrumentation phase, highlighting a certain fragility 
    of the tool to instrument any type of binary.

    Another important limitation that must be mentioned is that the presented approach is not thread-safe.
    Infact, the rax value is saved in a global memory area that is not thread safe.
    For this reason, if two parallel threads reach a \gls{SHM} update block at the same time, the value inside of save 
    rax register will be overwritten, causing one of the two threads to restore an incorrect value of rax.
    The same problem happens with previous \gls{BB} ID saved.
    Supposing that, at a certain point, four different threads are active, each thread will have, in general, a different 
    previous \gls{BB} ID.

    Another limitation worth mentioning is that the developed tool works only with \gls{ELF} files, lacking 
    compatibility with the Windows PE executable format.

\section{Future Works}
    The main work to be done in the future consists in making this type of solution more resilient to binary corruption.
    The approach of instrumenting the \gls{SHM} update directly inline within each \gls{BB} has proven effective, 
    given the speedup resulting from avoiding a context switch at each \gls{BB} (which would otherwise be caused 
    by invoking a function) [Reference to the experiments section].
    For this reason, a possible direction for future work could be to apply current state-of-the-art techniques 
    to this approach.
    For example, using the \gls{SHM} update rule presented in this thesis, several hash collisions occur in the 
    \gls{SHM}, causing the fuzzer to misinterpret coverage, as demonstrated by Shuitao Gan et al. \cite{CollAFL}.
    In fact, that paper presents a new approach to \gls{SHM} updates, which was later added as a supported 
    instrumentation method in the AFL++ framework. Implementing a collision-free coverage algorithm as 
    \gls{SHM} update rule could be a possible improvement over the proposed approach.

    Furthermore, implementing a sanitization mechanism could be another research direction building on the work 
    done in this thesis, for example by implementing an Address Sanitizer \cite{serebryany2012addresssanitizer} 
    or a Thread Sanitizer \cite{schilling2024binary}.
